% Generated draft: requires mathematical and source review.
\documentclass[11pt]{article}
\usepackage[margin=25mm]{geometry}
\usepackage{fontspec}
\setmainfont{Latin Modern Roman}
\usepackage{amsmath,amssymb,mathtools,unicode-math}
\setmathfont{Latin Modern Math}
\usepackage{enumitem,needspace,xurl,hyperref}
\setlength{\parindent}{0pt}
\setlength{\parskip}{5pt}
\setlength{\emergencystretch}{3em}
\newcommand{\N}{\mathbb{N}}
\newcommand{\Z}{\mathbb{Z}}
\newcommand{\Q}{\mathbb{Q}}
\newcommand{\R}{\mathbb{R}}
\newcommand{\C}{\mathbb{C}}
\newcommand{\F}{\mathbb{F}}
\newcommand{\A}{\mathbb{A}}
\newcommand{\PP}{\mathbb{P}}
\newcommand{\E}{\mathbb{E}}
\newcommand{\eps}{\varepsilon}
\newcommand{\dd}{\,\mathrm{d}}
\begin{document}
\section*{UnsolvedMath ID: 30004895}
\textbf{Source problem code:} catalogue entry
\subsection{Definitions and mathematical statement}
The supplied record concerns a distinguished presentation of a string cone by tropicalized inequalities, but does not specify the string-cone data, coordinates, or normalized defining forms. The following states the proposed criterion conditionally on those choices being supplied.

Let \(n,m\geq 1\), and let the designated defining linear forms, in the designated coordinates, be
\[
\ell_j(x)=\sum_{k=1}^{n}a_{jk}x_k\qquad (1\leq j\leq m).
\]
Using the convention that the inequalities are written with nonnegative right-hand orientation, let the associated string cone be
\[
C=\{x\in\R^n:\ell_j(x)\geq 0\text{ for every }1\leq j\leq m\}.
\]
This indexed presentation is \emph{irredundant} if removal of any one inequality strictly enlarges the feasible set; equivalently,
\[
\forall j\in\{1,\ldots,m\}\ \exists x\in\R^n:
\quad \ell_j(x)<0
\quad\text{and}\quad
\ell_i(x)\geq 0\ \text{for every }i\ne j.
\]
The proposed equivalence, for every presentation belonging to the intended source-specific class, is
\[
\bigl[\text{the designated presentation of }C\text{ is irredundant}\bigr]
\quad\Longleftrightarrow\quad
\bigl[\forall j\in\{1,\ldots,m\}\ \forall k\in\{1,\ldots,n\},\ a_{jk}\in\{-1,0,1\}\bigr].
\]
\emph{Unresolved qualifications.} The class of string cones and the distinguished presentation to which this quantifier applies are absent from the supplied excerpt. In particular, the Lie-theoretic data, any reduced-word choices, coordinate conventions, normalization of forms, and treatment of repeated inequalities must be recovered from the cited sources. The coefficient condition depends on normalization, whereas irredundancy is invariant under positive rescaling of individual forms. Thus the statement cannot be interpreted as a criterion for arbitrary linear-inequality presentations. The original excerpt refers to an omitted ``above theorem,'' so the displayed equivalence is a conditional reconstruction from the supplied clean statement, not a fully recovered source formulation.
\subsection{Short English statement}
The conjecture asks whether the designated tropical inequalities for a string cone are irredundant exactly when all their coefficients lie in \(\{-1,0,1\}\). The supplied record omits the precise family of cones and normalized inequalities, so those choices remain unresolved. Its open status and claim that the coefficient condition is sufficient are supplied dated assessments, not independently verified facts.
\subsection{Sources}
\begin{itemize}
\item Local registry: \texttt{data/unsolvedmath/problems.json}, permanent UnsolvedMath ID \(30004895\), source problem code \texttt{catalogue entry}, title ``Coefficient Criterion for Irredundant String-Cone Inequalities.'' The supplied record labels the problem open and its statement \texttt{corrected\_verified}; these are registry labels.
\item Bea Schumann, ``String cones and cluster varieties,'' joint work with Gleb Koshevoy, contribution identified in the supplied excerpt at p.~2883 of \emph{Enveloping Algebras and Geometric Representation Theory}. The supplied report citation is \emph{Oberwolfach Reports} (2022), pp.~2827--2891, DOI \texttt{10.4171/owr/2021/52}: \url{https://doi.org/10.4171/owr/2021/52}. The registry gives proposed year 2021. The original excerpt conjectures that an omitted preceding theorem is an equivalence.
\item Supplied literature reference: Koshevoy--Schumann, \emph{Journal of Algebraic Combinatorics}, \url{https://doi.org/10.1007/s10801-022-01144-z}.
\item Supplied literature reference: ``Oberwolfach Report 52, string cones,'' \url{https://doi.org/10.4171/owr/2021/52}.
\item Literature assessment dated exactly \texttt{2026-08-22}: the supplied curation describes the equivalence as open and states that the primary paper proves only the sufficient direction. This is supplied dated curation derived from the OWR material, not an independently verified literature assessment; time-sensitive status and the cited claims require independent verification.
\item The registry states that the official report was checked directly and that the clean statement restores definitions and preserves the mathematical target. That verification claim is supplied provenance; no source visit or independent verification is asserted here.
\end{itemize}
\end{document}
